\chapter[Karmaşıklık ve Big O · \textit{Temel}]{Karmaşıklık ve Big O}
\chaptermark{Karmaşıklık ve Big O}

Yazdığımız bir programın doğru çalışması tek başına yeterli değildir; programın verilen zaman ve bellek sınırları içinde de sonuç üretmesi gerekir. TÜBİTAK 1.~Aşama sınavında ve programlama yarışmalarında problemler genellikle katı süre sınırlarına (çoğunlukla 1 saniye) ve bellek sınırlarına (örneğin 256 MB) sahiptir. Aynı problemi çözen iki farklı algoritmadan biri sonucu saniyenin binde birinde bulurken, diğeri kabul edilebilir sürelerin çok ötesinde çalışabilir.

Algoritmaların çalışma hızını ve bellek tüketimini; donanımdan, işletim sisteminden ya da programlama dilinden bağımsız biçimde ölçmek temel bir gerekliliktir. Bu doğrultuda önce temel işlem adımlarını saymayı, ardından bu davranışı matematiksel olarak ifade eden Big O notasyonunu ve yarışmalardaki pratik sınırları ele alacağız.

\section{İşlem Sayısı}

Bir algoritmanın ne kadar hızlı çalıştığını saniye cinsinden ölçmek ilk başta mantıklı görünebilir. Ancak aynı programı eski bir dizüstü bilgisayarda çalıştırdığımızda 5 saniye, modern bir sunucuda çalıştırdığımızda ise 0.5 saniye sürebilir. Dahası, arka planda müzik çalması veya işletim sisteminin bir güncelleme yapması bile ölçülen süreyi değiştirebilir. 

Bize donanımdan bağımsız, evrensel bir ölçü gerekir: \textbf{Algoritmanın girdi boyutuna bağlı olarak yürüttüğü temel işlem sayısı}.

\subsection{Temel İşlem Nedir?}

Bir bilgisayarın işlemcisinin sabit sürede gerçekleştirebildiği en küçük adımlara \emph{temel işlem} diyoruz. Bunlar:
\begin{itemize}
    \item İki sayıyı toplamak, çıkarmak, çarpmak veya bölmek ($a + b$, $x \times y$),
    \item İki değeri karşılaştırmak ($a < b$, $x == y$),
    \item Bir değişkene değer atamak ($x \gets 5$),
    \item Dizinin belirli bir indisine erişmek ($A[i]$).
\end{itemize}

Şimdi somut bir problem üzerinden adım sayısını tam olarak hesaplayalım.

\begin{example}
$n$ elemanlı bir tam sayı dizisindeki en büyük elemanı bulmak istiyoruz. Aşağıdaki algoritmanın toplam işlem sayısını $n$ cinsinden belirleyelim.

\begin{lstlisting}[language=CTemel]
int en_buyuk = A[0];
for (int i = 1; i < n; i++) {
    if (A[i] > en_buyuk) {
        en_buyuk = A[i];
    }
}
\end{lstlisting}
\end{example}

\begin{proof}[Çözüm]
Adımları tek tek sayalım:
\begin{enumerate}
    \item \texttt{int en\_buyuk = A[0];}: 1 dizi erişimi ve 1 atama (2 işlem).
    \item Döngü başlangıcı: \texttt{int i = 1;} (1 atama).
    \item Döngü koşulu: \texttt{i < n} kontrolü her adımda yapılır. Döngü $n-1$ kez döner, son kontrolle birlikte $n$ kez karşılaştırma yapılır.
    \item İndis artırımı: \texttt{i++} işlemi $n-1$ kez yapılır.
    \item Karşılaştırma: \texttt{A[i] > en\_buyuk} koşulu her döngü adımında 1 dizi erişimi ve 1 karşılaştırma olmak üzere 2 işlem gerektirir, toplam $2(n-1)$ işlem.
    \item Güncelleme: \texttt{en\_buyuk = A[i]} satırı dizinin içeriğine bağlı olarak en iyi durumda 0 kez, en kötü durumda (dizi kesin artan ise) $n-1$ kez çalışır. Her çalışmada 1 dizi erişimi ve 1 atama (2 işlem) vardır.
\end{enumerate}

En kötü durumda (worst-case) toplam işlem sayısı:
\[ T(n) = 2 + 1 + n + (n-1) + 2(n-1) + 2(n-1) = 6n - 2 \]

En iyi durumda (best-case, ilk eleman en büyükse):
\[ T(n) = 2 + 1 + n + (n-1) + 2(n-1) + 0 = 4n \]

Görüldüğü üzere girdi boyutu $n$ iki katına çıktığında, yürütülen işlem sayısı da yaklaşık iki katına çıkmaktadır. Katsayılar (4 veya 6) donanıma veya derleyici optimizasyonuna göre değişebilir; ancak büyümenin $n$ ile \emph{doğrusal} olduğu gerçeği sabittir.
\end{proof}

\subsection{İç İçe Döngüler ve Büyüme Hızı}

Şimdi iki farklı algoritmayı karşılaştıralım. Amacımız $1$'den $n$'e kadar olan sayıların toplamını hesaplamak olsun.

\textbf{1. Yöntem (Tek Döngü):}
\begin{lstlisting}[language=CTemel]
long long toplam = 0;
for (int i = 1; i <= n; i++) {
    toplam += i;
}
\end{lstlisting}
Burada döngü $n$ defa döner. Her adımda 1 toplama ve 1 atama yapılır. Yaklaşık işlem sayısı $c_1 \cdot n$ kadardır (burada $c_1$ küçük bir sabittir).

\textbf{2. Yöntem (Matematiksel Formül):}
\begin{lstlisting}[language=CTemel]
long long toplam = (long long)n * (n + 1) / 2;
\end{lstlisting}
Burada $n = 10$ olsa da, $n = 10^9$ olsa da sadece 1 toplama, 1 çarpma ve 1 bölme yapılır. İşlem sayısı girdi boyutundan tamamen bağımsızdır; sabit bir $c_2$ sayısıdır.

\textbf{3. Yöntem (Gereksiz İç İçe Döngü):}
Aynı toplamı, $i$ sayısını $1$'den $i$'ye kadar $1$'leri toplayarak hesaplayan yapay bir algoritma düşünelim:
\begin{lstlisting}[language=CTemel]
long long toplam = 0;
for (int i = 1; i <= n; i++) {
    for (int j = 1; j <= i; j++) {
        toplam += 1;
    }
}
\end{lstlisting}
İçteki döngü $i = 1$ iken 1 kez, $i = 2$ iken 2 kez, \dots, $i = n$ iken $n$ kez çalışır. Toplam adım sayısı:
\[ 1 + 2 + 3 + \dots + n = \frac{n(n+1)}{2} = \frac{1}{2}n^2 + \frac{1}{2}n \]

Girdi boyutu büyüdükçe bu üç algoritmanın adım sayılarını kıyaslayalım:

\begin{center}
\begin{tabular}{|r|r|r|r|}
\hline
$n$ & 2. Yöntem ($c_2$) & 1. Yöntem ($2n$) & 3. Yöntem ($\frac{1}{2}n^2 + \frac{1}{2}n$) \\
\hline
$10$ & 3 & 20 & 55 \\
$1\,000$ & 3 & $2\,000$ & $500\,500$ \\
$100\,000$ & 3 & $200\,000$ & $5\,000\,050\,000$ ($5 \times 10^9$) \\
$1\,000\,000$ & 3 & $2\,000\,000$ & $500\,000\,500\,000$ ($5 \times 10^{11}$) \\
\hline
\end{tabular}
\end{center}

$n = 100\,000$ olduğunda 1. yöntem modern bir bilgisayarda birkaç milisaniyede biterken, 3. yöntem saniyeler sürer. $n = 1\,000\,000$ olduğunda 3. yöntemin bitmesi dakikalar alacaktır.

Buradan çıkaracağımız temel sonuç şudur: \textbf{Girdi $n$ büyüdükçe, fonksiyonun en yüksek dereceli terimi tüm davranışı belirler.} $T(n) = \frac{1}{2}n^2 + \frac{1}{2}n$ ifadesinde $n = 10^6$ için $\frac{1}{2}n^2 = 5 \times 10^{11}$ iken, $\frac{1}{2}n = 5 \times 10^5$'tir. Küçük dereceli terim, toplamın yanında ihmal edilebilir düzeydedir. Benzer biçimde öndeki $\frac{1}{2}$ katsayısı da $n^2$'nin büyüme hızının yanında belirleyici rolünü kaybeder.

Bu gözlem bizi algoritmik karmaşıklığın standart diline götürür: Big O notasyonu.

\section{Big O Notasyonu}

\subsection{Motivasyon ve Biçimsel Tanım}

Bir algoritmanın işlem sayısı tam olarak $T(n) = 3n^2 + 14n + 120$ olsun. Başka bir algoritmanınki ise $T'(n) = n^2 + 200n + 5$ olsun. İkisi de temelde $n^2$ ile ölçeklenir; yani girdi 10 katına çıktığında işlem sayısı kabaca 100 katına çıkar. Analiz yaparken ayrıntılardaki $14n$ veya sabit $120$ gibi küçük terimlerle ve baştaki katsayılarla uğraşmak istemeyiz. Algoritmanın \emph{asemptotik} (yani $n$ sonsuza giderkenki) üst sınırını belirlemek isteriz.

\begin{definition}[Big O (Büyük O)]
$f(n)$ ve $g(n)$, pozitif tam sayılardan pozitif reel sayılara tanımlı iki fonksiyon olsun. Eğer her $n \ge n_0$ için
\[ f(n) \le c \cdot g(n) \]
olacak biçimde pozitif reel bir $c$ sabiti ve pozitif bir $n_0$ tam sayısı varsa,
\[ f(n) = O(g(n)) \]
yazılır ve ``$f(n)$, $g(n)$ mertebesindedir'' denir.
\end{definition}

Bu tanım şunu söyler: Yeterince büyük $n$ değerleri için ($n \ge n_0$), $f(n)$ fonksiyonu $g(n)$'in sabit bir katını asla geçemez. $g(n)$, $f(n)$ için bir \textbf{asemptotik üst sınırdır}.

\begin{theorem}
$P(n) = a_k n^k + a_{k-1} n^{k-1} + \dots + a_1 n + a_0$ derecesi $k$ olan bir polinom ve $a_k > 0$ olsun. Bu durumda $P(n) = O(n^k)$'dir.
\end{theorem}

\begin{proof}
Her $n \ge 1$ için üçgen eşitsizliğini kullanarak:
\[ |P(n)| \le |a_k| n^k + |a_{k-1}| n^{k-1} + \dots + |a_0| \]
$n \ge 1$ olduğundan her $i \le k$ için $n^i \le n^k$ yazabiliriz. Dolayısıyla:
\[ |P(n)| \le (|a_k| + |a_{k-1}| + \dots + |a_0|) n^k \]
Burada $c = |a_k| + |a_{k-1}| + \dots + |a_0|$ ve $n_0 = 1$ seçilirse tanımdaki eşitsizlik sağlanır. Demek ki $P(n) = O(n^k)$'dir.
\end{proof}

Bu teorem sayesinde karmaşık polinomlarla tek tek uğraşmak gerekmez: En büyük dereceli terimi alır, katsayısını 1 yaparız. Örneğin:
\[ 5n^3 + 100n^2 - 7n + 42 = O(n^3) \]

\subsection{Temel Karmaşıklık Sınıfları ve Sezgileri}

Olimpiyatlarda en sık karşılaşacağımız karmaşıklık sınıflarını küçükten büyüğe sıralayalım:

\begin{multline*}
O(1) < O(\log n) < O(\sqrt{n}) < O(n) < O(n \log n) \\
< O(n^2) < O(n^3) < O(2^n) < O(n!)
\end{multline*}

Şimdi bu sınıfların sezgisel anlamlarını inceleyelim.

\subsubsection{1. Sabit Zaman: $O(1)$}
İşlem sayısı girdinin büyüklüğünden bağımsızdır. Girdi ister 5 ister 5 milyar olsun, işlem sabit sayıda adımda tamamlanır.
\begin{itemize}
    \item Bir dizinin $k$. elemanına erişmek: \texttt{A[k]}
    \item İki sayının tek mi çift mi olduğunu kontrol etmek: \texttt{x \% 2 == 0}
    \item Aritmetik bir formülü hesaplamak: \texttt{n * (n + 1) / 2}
\end{itemize}

\subsubsection{2. Logaritmik Zaman: $O(\log n)$}
Bir algoritma her adımda arama uzayını veya girdiyi sabit bir oranda (genellikle yarı yarıya) küçültüyorsa, o algoritmanın karmaşıklığı logaritmiktir. 

Örnek olarak bir sayıyı sürekli 2'ye böldüğümüzü düşünelim:
\begin{lstlisting}[language=CTemel]
int sayac = 0;
while (n > 1) {
    n = n / 2;
    sayac++;
}
\end{lstlisting}
$n = 16$ için değerler: $16 \to 8 \to 4 \to 2 \to 1$ (4 adım). 
Genel olarak $2^k \le n$ ise döngü yaklaşık $\lfloor \log_2 n \rfloor$ adım sürer. Big O gösteriminde logaritmanın tabanı yazılmaz; çünkü taban değişimi yalnızca sabit bir çarpan getirir: $\log_a n = \frac{\log_b n}{\log_b a} = c \cdot \log_b n$. Dolayısıyla $O(\log_2 n) = O(\log_{10} n) = O(\log n)$'dir.

$O(\log n)$ algoritmaları son derece yavaş büyür. $n = 10^{18}$ gibi çok büyük bir girdi için bile $\log_2(10^{18}) \approx 60$ adımdır.

\subsubsection{3. Doğrusal Zaman: $O(n)$}
Girdideki her elemana sabit sayıda (örneğin 1 veya 2 kez) bakılıyorsa algoritma doğrusaldır.
\begin{itemize}
    \item $n$ elemanlı bir dizinin toplamını bulmak.
    \item Sırasız bir dizide belirli bir değeri aramak (lineer arama).
\end{itemize}

\subsubsection{4. Doğrusal-Logaritmik Zaman: $O(n \log n)$}
Genellikle bir problemi logaritmik derinlikte parçalara ayırıp her seviyede $O(n)$ iş yapan ``böl ve yönet'' algoritmalarında karşımıza çıkar. Yaygın sıralama algoritmaları (Merge Sort, Hızlı Sıralama) bu sınıftadır. $n = 10^6$ için $n \log_2 n \approx 2 \times 10^7$ işlem demektir ve bu olimpik süre sınırları için oldukça güvenli bir çalışma hacmidir.

\subsubsection{5. Karesel Zaman: $O(n^2)$}
Genellikle iç içe iki döngünün tüm çiftleri taradığı durumlarda ortaya çıkar.
\begin{lstlisting}[language=CTemel]
for (int i = 0; i < n; i++) {
    for (int j = 0; j < n; j++) {
        // Sabit sayida islem O(1)
    }
}
\end{lstlisting}
Tüm $(i, j)$ ikililerinin sayısı $n \times n = n^2$'dir. $n = 1\,000$ için $n^2 = 10^6$ (hızlı), fakat $n = 100\,000$ için $n^2 = 10^{10}$ olur ve süre sınırını aşar.

\subsubsection{6. Üstel ve Faktöriyel Zaman: $O(2^n)$ ve $O(n!)$}
Bölüm 4 ve 5'te gördüğümüz sayma kurallarıyla bağlantılı olarak, bir kümenin tüm alt kümelerini gezmek $2^n$, tüm permütasyonlarını gezmek ise $n!$ adım gerektirir. Bu algoritmalar yalnızca çok küçük $n$ değerleri ($n \le 20$ veya $n \le 11$) için pratik sürede çalışabilir.

\subsection{Çözümlü Örnekler: Karmaşıklık Nasıl Hesaplanır?}

\begin{example}
Aşağıdaki kod parçasının zaman karmaşıklığını Big O cinsinden bulunuz.

\begin{lstlisting}[language=CTemel]
for (int i = 1; i <= n; i *= 2) {
    for (int j = 1; j <= n; j++) {
        toplam++;
    }
}
\end{lstlisting}
\end{example}

\begin{proof}[Çözüm]
İç ve dış döngülerin adım sayılarını birbirinden bağımsız biçimde inceleyebiliriz:
\begin{itemize}
    \item Dıştaki döngü: $i$ değişkeni $1, 2, 4, 8, \dots$ şeklinde ikiye katlanarak ilerliyor. Dolayısıyla $i \le n$ koşulu sağlandığı sürece döngü adımı sayısı $k$, $2^{k-1} \le n$ eşitsizliğine uyar. Yani dış döngü tam olarak $\lfloor \log_2 n \rfloor + 1$ kez döner.
    \item İçteki döngü: Dış döngünün her bir adımında, $j$ değişkeni $1$'den $n$'e kadar birer birer artar. Yani iç döngü her zaman $n$ kez çalışır.
\end{itemize}
Toplam işlem sayısı iç ve dış döngü adımlarının çarpımıdır:
\[ (\lfloor \log_2 n \rfloor + 1) \times n = O(n \log n) \]
\end{proof}

\begin{example}
Aşağıdaki kod parçasının zaman karmaşıklığını belirleyiniz.

\begin{lstlisting}[language=CTemel]
for (int i = 1; i <= n; i++) {
    for (int j = 1; j <= n; j += i) {
        toplam++;
    }
}
\end{lstlisting}
\end{example}

\begin{proof}[Çözüm]
İçteki döngünün adım sayısı $i$'ye bağlı olduğundan sınırları doğrudan çarpmak yerine her $i$ değeri için adımları tek tek toplayalım:

\begin{itemize}
    \item $i = 1$ için: $j$ birer birer artar, döngü $\frac{n}{1}$ kez çalışır.
    \item $i = 2$ için: $j$ ikişer ikişer artar, döngü $\frac{n}{2}$ kez çalışır.
    \item $i = 3$ için: $j$ üçer üçer artar, döngü $\frac{n}{3}$ kez çalışır.
    \item \dots
    \item $i = n$ için: $j$ $n$'er artar, döngü $\frac{n}{n} = 1$ kez çalışır.
\end{itemize}

Toplam adım sayısı:
\[ S = \frac{n}{1} + \frac{n}{2} + \frac{n}{3} + \dots + \frac{n}{n} = n \left( 1 + \frac{1}{2} + \frac{1}{3} + \dots + \frac{1}{n} \right) \]

Parantez içindeki ifade harmonik seridir ($H_n$). Bilindiği üzere $H_n = \ln n + \gamma + O(1/n)$ olup büyüme hızı $O(\log n)$'dir. 

Dolayısıyla toplam adım sayısı:
\[ S = n \cdot O(\log n) = O(n \log n) \]
Bu analiz, Bölüm 12'de göreceğimiz Eratosthenes kalburunun karmaşıklık hesabının da temelini oluşturur.
\end{proof}

\begin{example}
Aşağıdaki fonksiyonun zaman karmaşıklığını bulunuz.

\begin{lstlisting}[language=CTemel]
void islem(int n) {
    while (n > 0) {
        for (int i = 0; i < n; i++) {
            printf("*");
        }
        n /= 2;
    }
}
\end{lstlisting}
\end{example}

\begin{proof}[Çözüm]
İlk bakışta iç içe bir yapı görüp doğrudan $O(n \log n)$ demek sık yapılan bir hatadır. Adımları toplayarak görelim:
\begin{itemize}
    \item 1. adımda $n$ yıldız basılır.
    \item 2. adımda $n / 2$ yıldız basılır.
    \item 3. adımda $n / 4$ yıldız basılır.
    \item \dots
\end{itemize}
Toplam işlem sayısı:
\[ T(n) = n + \frac{n}{2} + \frac{n}{4} + \frac{n}{8} + \dots + 1 \]
$n$ parantezine alırsak:
\[ T(n) = n \left( 1 + \frac{1}{2} + \frac{1}{4} + \frac{1}{8} + \dots \right) \]
Sonsuz geometrik serinin toplamı $\sum_{i=0}^{\infty} \left(\frac{1}{2}\right)^i = \frac{1}{1 - 1/2} = 2$'dir.
Dolayısıyla:
\[ T(n) < 2n \]
Yani toplam işlem sayısı $2n$'i asla geçemez. O halde bu algoritmanın karmaşıklığı $O(n \log n)$ değil, \textbf{$O(n)$}'dir.
\end{proof}

\begin{example}[Döngü Gövdesindeki Gizli Maliyetler]
Aşağıdaki kod parçasının karmaşıklığını inceleyelim:

\begin{lstlisting}
result = empty list
FOR EACH element IN A DO
    IF element NOT IN result THEN
        APPEND element TO result
    END IF
END FOR\end{lstlisting}
\end{example}

\begin{proof}[Çözüm]
Tek bir döngü görüldüğünde kodun $O(n)$ olduğunu varsaymak sık yapılan bir yanılgıdır.

İçerideki arama adımı (\texttt{element NOT IN result}), bir listenin içinde baştan sona eleman arar. \texttt{result} listesinde $k$ eleman varken bu arama işlemi $O(k)$ sürer. Dizi tekil elemanlardan oluşuyorsa $k$, $0$'dan $n-1$'e kadar büyür. 
Toplam maliyet:
\[ 0 + 1 + 2 + \dots + (n-1) = \frac{n(n-1)}{2} = O(n^2) \]
Programlama dillerinin yerleşik fonksiyonlarını (örneğin dizide arama, dilimleme ya da kopyalama) kullanırken o işlemlerin arka plandaki maliyetini mutlaka hesaba katmalıyız.
\end{proof}

\section{Pratik Limitler}

Teorik Big O analizini somut yarışma koşullarına uyarlayalım. Karşımızda girdi boyutu $n$ olarak verilmiş ve süre sınırı 1.0 saniye olan bir problem olduğunda, hangi karmaşıklıktaki bir algoritmayı tasarlamalıyız?

\subsection{Temel Eşik: $10^8$ İşlem / Saniye}

\begin{definition}[$10^8$ Kuralı]
Modern yarışma sunucularında (Codeforces, TÜBİTAK sınav sistemi, AtCoder vb.) standart bir C/C++ programı \textbf{1 saniyede yaklaşık $10^8$ (100 milyon) basit işlem} gerçekleştirir.
\end{definition}

Bu kural mutlak bir sınır olmamakla birlikte yürütülen işlemlerin türüne göre esneklik gösterir:
\begin{itemize}
    \item Basit toplama, çıkarma, bit işlemleri (bitwise) ve dizi erişimleri çok hızlıdır; 1 saniyede $2 \times 10^8$ hatta $4 \times 10^8$ adet çalışabilir.
    \item Çarpma nispeten hızlıdır; bölme (\texttt{/}) ve mod alma (\texttt{\%}) işlemleri ise işlemci için daha maliyetli olup 1 saniyede ancak $5 \times 10^7$ civarında yürütülebilir.
    \item Kayan noktalı sayı fonksiyonları (\texttt{sqrt}, \texttt{sin}, \texttt{pow}) belirgin biçimde daha yavaştır.
\end{itemize}

Bununla birlikte, kaba bir ön tahmin için $10^8$ eşiği en güvenilir başlangıç noktasıdır.

\subsection{Girdi Boyutuna ($n$) Göre Algoritma Tahmini}

TÜBİTAK 1.~Aşama sınavında problem metnindeki $n$ sınırlarına bakarak hedeflenen algoritmanın karmaşıklığı doğrudan tahmin edilebilir. Aşağıdaki tablo bu kararda temel rehberdir:

\begin{center}
\begin{tabular}{|l|p{4.2cm}|p{4.6cm}|}
\hline
\textbf{Girdi Boyutu ($n$)} & \textbf{Beklenen Karmaşıklık} & \textbf{Tipik Algoritma / Yaklaşım} \\
\hline
$n \le 10$ & $O(n!)$ & Tüm permütasyonları deneme \\
$n \le 20$ & $O(2^n)$ veya $O(2^n \cdot n)$ & Tüm alt kümeleri gezme, bitmask DP \\
$n \le 500$ & $O(n^3)$ & Floyd-Warshall, 3 iç içe döngü \\
$n \le 5\,000$ & $O(n^2)$ & Kaba kuvvet ikilileri tarama, basit DP \\
$n \le 200\,000$ & $O(n \log n)$ veya $O(n)$ & Sıralama, ikili arama, segment ağacı \\
$n \le 10^7$ & $O(n)$ & Basit tek geçişli döngü, kalbur \\
$n > 10^9$ & $O(\log n)$ veya $O(1)$ & Matematiksel formül, hızlı üs alma \\
\hline
\end{tabular}
\end{center}

\begin{example}
Bir problemde $n = 100\,000$ tam sayı veriliyor. Tüm $(i, j)$ indis çiftleri ($1 \le i < j \le n$) için $|A[i] - A[j]|$ farklarının en küçüğünü bulmamız isteniyor. Hangi yaklaşımı seçmeliyiz?
\end{example}

\begin{proof}[Çözüm]
\textbf{1. Yaklaşım (Kaba Kuvvet):} Tüm çiftleri iç içe iki döngüyle deneriz. 
Bölüm 5'te gördüğümüz kombinasyon formülüyle, tüm çiftlerin sayısı $\binom{n}{2} = \frac{n(n-1)}{2}$'dir. 
$n = 10^5$ için:
\[ \frac{10^5 \times (10^5 - 1)}{2} \approx \frac{10^{10}}{2} = 5 \times 10^9 \text{ işlem} \]
$10^8$ kuralına göre bu işlem yaklaşık $\frac{5 \times 10^9}{10^8} = 50$ saniye sürer. Süre sınırı 1 saniye olduğundan bu yöntem \textbf{Zaman Sınırı Aşımı} (Time Limit Exceeded - TLE) verecektir.

\textbf{2. Yaklaşım (Sıralama):} Sayıları önce küçükten büyüğe sıralarız. Birbirine en yakın iki sayı, sıralı dizide kesinlikle ardışık olmak zorundadır.
\begin{enumerate}
    \item Diziyi sırala: $O(n \log n)$
    \item Yan yana duran elemanların farkına tek bir döngüde bak: $O(n)$
\end{enumerate}
Toplam karmaşıklık $O(n \log n)$ olur. $n = 10^5$ için:
\[ n \log_2 n \approx 10^5 \times 17 \approx 1.7 \times 10^6 \text{ işlem} \]
Bu da yaklaşık $0.02$ saniye sürer ve süre sınırının oldukça altında kalır.
\end{proof}

\subsection{Alan (Bellek) Karmaşıklığı}

Algoritmalar yalnızca zaman değil, hafıza da tüketir. Yarışmalarda tipik bellek sınırı 128 MB ile 512 MB arasındadır (en yaygın olanı 256 MB).

Bellek analizinde temel veri tiplerinin kapladığı yerleri bilmek gerekir:
\begin{itemize}
    \item \texttt{bool} / \texttt{char}: 1 bayt
    \item \texttt{int}: 4 bayt
    \item \texttt{long long} / \texttt{double}: 8 bayt
\end{itemize}

Bellek miktarını megabayta çevirmek için:
\[ 1 \text{ KB} = 1024 \text{ bayt}, \quad 1 \text{ MB} = 1024 \text{ KB} \approx 10^6 \text{ bayt} \]

Pratik kural olarak:
\begin{itemize}
    \item $10^7$ elemanlı bir \texttt{int} dizisi: $10^7 \times 4 \text{ bayt} \approx 40 \text{ MB}$ yer tutar.
    \item $10^8$ elemanlı bir \texttt{int} dizisi: $\approx 400 \text{ MB}$ yer tutar (ve 256 MB sınırında hafıza aşımı verir).
    \item $5000 \times 5000$ boyutunda iki boyutlu bir \texttt{int} matrisi: $2.5 \times 10^7 \times 4 \approx 100 \text{ MB}$ yer tutar.
\end{itemize}

\begin{example}[Alan-Zaman Takası (Space-Time Tradeoff)]
Algoritma tasarımında sıklıkla ek bellek kullanılarak hız kazanılır veya tam tersine bellek tasarrufu adına işlem süresi artırılır.

Bir dizideki sayıların daha önce görülüp görülmediğini denetlemek için:
\begin{itemize}
    \item Ekstra bellek kullanmazsak: Her yeni eleman için dizinin başına bakarız $\to O(n^2)$ zaman, $O(1)$ ek bellek.
    \item Bir frekans dizisi (veya doğruluk tablosu) tutarsak: Sayıyı gördüğümüzde \texttt{var\_mi[x] = true} yaparız $\to O(n)$ zaman, sayının alabileceği en büyük değer ($M$) kadar $O(M)$ ek bellek.
\end{itemize}
\end{example}

\begin{example}[Rekürsiyon Yığını (Call Stack)]
Yalnızca tanımladığınız diziler değil, çağırdığınız rekürsif (rekürsif) fonksiyonlar da bellek tüketir. Her fonksiyon çağrısı işletim sisteminin çağrı yığını (call stack) bölgesinde yer tutar.

Örneğin:
\begin{lstlisting}[language=CTemel]
void dfs(int u) {
    for (int v : komsular[u]) {
        if (!ziyaret[v]) dfs(v);
    }
}
\end{lstlisting}
Graf bir zincir biçimindeyse ve $n = 200\,000$ ise, \texttt{dfs} fonksiyonu iç içe $200\,000$ kez çağrılır. Her çağrı yığında yaklaşık 32--64 bayt tuttuğunda toplam yığın boyutu $200\,000 \times 48 \approx 10 \text{ MB}$ seviyesine ulaşır. Pek çok yarışma platformunda ve Windows sistemlerinde varsayılan yığın limiti 8 MB civarındadır. Program genel bellek sınırını aşmamış olsa dahi \textbf{Stack Overflow} (Yığın Taşması) hatası alarak çöker.
\end{example}

\subsection{Bölüm Özeti ve Kontrol Listesi}

Bir problemi çözerken şu adımları alışkanlık haline getiriniz:
\begin{enumerate}
    \item Soru metnindeki $n$ sınırına bakarak hedef karmaşıklığı belirleyin ($n = 10^5$ ise $O(n^2)$ bir algoritma tasarlamak yerine doğrudan $O(n \log n)$ veya $O(n)$ çözümlere yönelin).
    \item Kodunuzdaki iç içe döngülerin sınırlarını ve yerleşik dil fonksiyonlarının maliyetlerini tek tek hesaba katın.
    \item İşlem sayınızın $10^8$ sınırını aşıp aşmadığını kaba bir hesapla test edin.
    \item Tanımladığınız dizilerin toplam boyutunun 256 MB limitini geçmediğinden emin olun.
\end{enumerate}
